\renewcommand{\Anno}{2005}

\renewcommand{\NumeroProblema}{2}

\renewcommand{\Testo}{La figura sottostante, in cui il lato di ogni
quadratino corrisponde a 100 metri, rappresenta un fiume ed un suo affluente.  Nella zona è possibile costruire strade lungo qualunque tracciato, con l'unico vincolo che gli eventuali ponti devono risultare perpendicolari alla direzione dei rispettivi corsi d'acqua.
	
\begin{center}
		\psset{unit=4.5mm}
		\pspicture(-8.5,-8.5)(8.5,8.5)
	
		\newrgbcolor{colorefiume}{0.6 0.8 1}
		\pspolygon*[linecolor=colorefiume]
		(-8,-8)(-7,-8)(8,7)(8,8)(7,8)(4,5)(-8,5)(-8,4)(3,4)(-8,-7)(-8,-8)
	
		\psline{-}(-8,-7)(3,4)(-8,4)
		\psline{-}(-8,5)(4,5)(7,8)
		\psline{-}(-7,-8)(8,7)
	
		\psdots[dotscale=1](-6,7)(4,-7)(-2,3)
		\uput[0](-6,7){$A$}
		\uput[225](-2,3){$B$}
		\uput[90](4,-7){$C$}
	
		\multido{\n=-8+1}{17}{
		\psline[linestyle=dotted,dotsep=1pt]{-}(-8,\n)(8,\n)
		\psline[linestyle=dotted,dotsep=1pt]{-}(\n,-8)(\n,8)}
	
		\endpspicture
\end{center}

	\begin{enumerate}
		\item Determinare la lunghezza minima possibile per una strada (eventuali ponti compresi) che unisce il punto $A$ ed il punto $B$.
	
		\item Determinare la lunghezza minima possibile per una strada (eventuali ponti compresi) che unisce il punto $A$ ed il punto $C$.
	\end{enumerate}}

\renewcommand{\Risposta}{\begin{enumerate}
		\item[(a)]  600 metri
	
		\item[(b)]  $(1600+100\sqrt{2})$ metri
	\end{enumerate}}

\renewcommand{\Aiutino}{\begin{enumerate}
	\renewcommand{\labelenumi}{(\alph{enumi})} 
	\item[(a)]  Supponendo di ``contrarre il fiume'', rimane una strada che unisce il punto $A$ con un punto $B'$ che si trova un quadretto al di sopra di $B$. Tale strada deve avere lunghezza minima.
	
	In alternativa, si può calcolare la lunghezza del percorso in funzione del punto in cui viene fatto il ponte, e determinarne il minimo studiando la funzione.

	\item[(b)]  In questo caso i fiumi da contrarre sono 2. Dopo averlo fatto, rimane una strada che unisce il punto $A$ con un punto $C'$ che si trova due quadretti sopra ed un quadretto a sinistra rispetto al punto $C$.
\end{enumerate}}

\renewcommand{\SoluzioneUfficiale}{\TitoloSoluzione \DomandaA È chiaro che
tra $A$ e $B$ almeno un ponte sull'affluente ci deve essere, ed è altrettanto chiaro che non conviene fare percorsi con più di un ponte. Ciò premesso, consideriamo una qualunque strada tra $A$ e $B$ (Figura 1). Poichè il ponte è perpendicolare al corso d'acqua, eliminando quest'ultimo, cioè traslando in su di un quadratino la sua ``sponda bassa'', i due estremi del ponte vanno a coincidere ed otteniamo un cammino tra $A$ e l'immagine $B'$ del punto $B$ (Figura 2).
\begin{center}
	\hfill
	\psset{unit=3ex}
		\pspicture(-1,-2)(7,6)
	
		\newrgbcolor{colorefiume}{0.6 0.8 1}
		
		\psframe*[linecolor=colorefiume](-1,2)(7,3)
	
		\psline{-}(-1,2)(7,2)
		\psline{-}(-1,3)(7,3)
		
		\pstGeonode(1,5){A}
		\pstGeonode(5,1){B}
	
		\pstGeonode[PointSymbol=none](3,3){P'}
		\pstGeonode[PointSymbol=none](3,2){Q'}
	
		\multips(-1,0)(1,0){9}{
		\psline[linestyle=dotted,dotsep=1pt]{-}(0,0)(0,6)}
		\multips(-1,0)(0,1){7}{
		\psline[linestyle=dotted,dotsep=1pt]{-}(0,0)(8,0)}
		
		\psline{->}(7.5,2)(7.5,3)
		\psline{->}(-1.5,2)(-1.5,3)
	
		\psset{linecolor=blue,linewidth=1.3\pslinewidth}
		\pstLineAB{A}{P'}
		\pstLineAB{Q'}{P'}
		\pstLineAB{Q'}{B}
		
		\rput{0}(3,-1){Figura 1}
		
	\endpspicture
	\hfill
	\psset{unit=3ex}
		\pspicture(-1,-2)(7,6)
	
		\newrgbcolor{colorefiume}{0.6 0.8 1}
		
	% 	\psframe*[linecolor=colorefiume](-1,2)(7,3)
	
	% 	\psline{-}(-1,2)(7,2)
		\psline{-}(-1,3)(7,3)
		
		\pstGeonode(1,5){A}
		\pstGeonode(5,1){B}
		\pstGeonode(5,2){B'}
	
		\pstGeonode[PointSymbol=none](3,3){P'}
		\pstGeonode[PointSymbol=none](3,2){Q'}
	
		\multips(-1,0)(1,0){9}{
		\psline[linestyle=dotted,dotsep=1pt]{-}(0,0)(0,6)}
		\multips(-1,0)(0,1){7}{
		\psline[linestyle=dotted,dotsep=1pt]{-}(0,0)(8,0)}
	
	% 	\pstLineAB[arrows=->]{B}{B'}
		\psset{linecolor=blue,linewidth=1.3\pslinewidth}
		\pstLineAB{A}{P'}
		\pstLineAB{B'}{P'}
		\pstLineAB[linestyle=dashed]{Q'}{P'}
		\pstLineAB[linestyle=dashed]{Q'}{B}
		
		\rput{0}(3,-1){Figura 2}
		
	\endpspicture
	\hfill\mbox{}
\end{center}

Per minimizzare la lunghezza totale non ci resta quindi che minimizzare la lunghezza di questo cammino fittizio tra $A$ e $B'$, il che si ottiene prendendo il segmento $AB'$, la cui lunghezza (per il teorema di Pitagora) è di 500 metri. Considerato che il ponte è lungo necessariamente 100 metri, avremo che la minima lunghezza per una strada che va da $A$ a $B$ è di 600 metri.

La figura 3 mostra la strada ottimale ricostruita a partire dal segmento $AB'$.
\begin{center}
	\psset{unit=3ex}
		\pspicture(-1,-2)(7,6)
	
		\newrgbcolor{colorefiume}{0.6 0.8 1}
		
		\psframe*[linecolor=colorefiume](-1,2)(7,3)
	
		\psline{-}(-1,2)(7,2)
		\psline{-}(-1,3)(7,3)
		
		\pstGeonode(1,5){A}
		\pstGeonode(5,1){B}
		\pstGeonode[PosAngle=45](5,2){B'}
	
		\pstGeonode[PointSymbol=none](-1,3){F}
		\pstGeonode[PointSymbol=none](7,3){F'}
		\pstGeonode[PointSymbol=none](3,3){P'}
		\pstGeonode[PointSymbol=none](3,2){Q'}
		
		\multips(-1,0)(1,0){9}{
		\psline[linestyle=dotted,dotsep=1pt]{-}(0,0)(0,6)}
		\multips(-1,0)(0,1){7}{
		\psline[linestyle=dotted,dotsep=1pt]{-}(0,0)(8,0)}
		
		\pstInterLL[PointSymbol=none]{A}{B'}{F}{F'}{P}
		\pstTranslation[PointSymbol=none]{P'}{Q'}{P}[Q]
		
		\psset{linecolor=blue,linewidth=1.3\pslinewidth}
		\pstLineAB[linestyle=dashed]{P}{B'}
		\pstLineAB{Q}{P}
		\pstLineAB{Q}{B}
		\pstLineAB{A}{P}
		
		\rput{0}(3,-1){Figura 3}
		
	\endpspicture
\end{center}

\DomandaB In questo caso è evidente che la strada migliore
dovrà avere 2 ponti, uno sull'affluente ed uno sul fiume. Ragioniamo ora come nel caso precedente: la contrazione del ponte sull'affluente sposta il punto $C$ di un quadratino verso l'alto, la contrazione del ponte sul fiume lo sposta di un altro quadratino verso l'alto e di uno verso sinistra.

Ci ritroviamo quindi un cammino che unisce $A$ con un punto $C'$ che da $A$ dista 900 metri in orizzontale e 1200 metri in verticale: per il teorema di Pitagora la distanza tra $A$ e $C'$ è dunque di 1500 metri.  Per avere il cammino di lunghezza minima dobbiamo aggiungere a questa 100 metri per il ponte sull'affluente e $100\sqrt{2}$ metri per il ponte sul fiume, per un totale di $(1600+100\sqrt{2})$ metri.

La figura 4 mostra il cammino ottimale ottenuto a partire dal segmento $AC'$ (i 3 tratti di strada ``terrestre'' sono tutti paralleli ad $AC'$).
\begin{center}
	\psset{unit=3ex}
		\pspicture(-8,-10)(8,8)
		
		\newrgbcolor{colorefiume}{0.6 0.8 1}
		\pspolygon*[linecolor=colorefiume]
		(-8,-8)(-7,-8)(8,7)(8,8)(7,8)(4,5)(-8,5)(-8,4)(3,4)(-8,-7)(-8,-8)
		
		\psline{-}(-8,-7)(3,4)(-8,4)
		\psline{-}(-8,5)(4,5)(7,8)
		\psline{-}(-7,-8)(8,7)
		
		\pstGeonode(-6,7){A}
		\pstGeonode(4,-7){C}
		\pstGeonode(3,-5){C'}

		\pstGeonode[PointSymbol=none](-8,5){F}
		\pstGeonode[PointSymbol=none](-7,5){F'}
		\pstGeonode[PointSymbol=none](7,8){G}
		\pstGeonode[PointSymbol=none](6,7){G'}
		\pstGeonode[PointSymbol=none](0,3){P'}
		\pstGeonode[PointSymbol=none](0,2){Q'}
		\pstGeonode[PointSymbol=none](1,1){R'}
			
		\pstInterLL[PointSymbol=none]{A}{C'}{F}{F'}{P}
		\pstTranslation[PointSymbol=none]{P'}{Q'}{P}[Q]
		\pstTranslation[PointSymbol=none]{A}{C'}{Q}[C'']
		\pstInterLL[PointSymbol=none]{Q}{C''}{G}{G'}{P''}
		\pstTranslation[PointSymbol=none]{Q'}{R'}{P''}[Q'']

		\multido{\n=-8+1}{17}{
		\psline[linestyle=dotted,dotsep=1pt]{-}(-8,\n)(8,\n)
		\psline[linestyle=dotted,dotsep=1pt]{-}(\n,-8)(\n,8)}
		
		\psset{linecolor=blue,linewidth=1.3\pslinewidth}
		\pstLineAB[linestyle=dashed]{P}{C'}
		\pstLineAB{A}{P}
		\pstLineAB{Q}{P}
		\pstLineAB{Q}{P''}
		\pstLineAB{P''}{Q''}
		\pstLineAB{Q''}{C}
		\rput{0}(0,-9){Figura 4}
		
	\endpspicture
\end{center}
}

\renewcommand{\SoluzioniAlternative}{\TitoloSoluzioneAlternativaA 
Mostriamo un approccio computazionale al punto (a). Come nella soluzione precedente osserviamo che possiamo limitarci alle strade con un solo ponte (sull'affluente) e che hanno andamento rettilineo sia prima sia dopo il ponte.  Con riferimento alla figura 5, in cui le lunghezze sono pensate in centinaia di metri, il problema si riduce pertanto a minimizzare la somma
$$AP+PQ+QB=\sqrt{4+x^{2}}+1+\sqrt{1+(4-x)^{2}},$$
in cui $x$ indica la distanza tra $P$ e la proiezione di $A$ sulla sponda più vicina dell'affluente.

\begin{center}
	\psset{unit=4ex}
		\pspicture(-1,-2)(7,6)
	
		\newrgbcolor{colorefiume}{0.6 0.8 1}
		
		\psframe*[linecolor=colorefiume](-1,2)(7,3)
	
		\psline{-}(-1,2)(7,2)
		\psline{-}(-1,3)(7,3)
		
		\pstGeonode(1,5){A}
		\pstGeonode[PosAngle=-45](5,1){B}
	
		\pstGeonode[PosAngle=45](2.3,3){P}
		\pstGeonode[PosAngle=-90](2.3,2){Q}
		\pstGeonode[PointName=none](1,3){H}
		\pstGeonode[PointName=none](5,2){K}
		
		\pcline[linestyle=none](H)(P)\aput[3pt]{0}{$x$}
		\pcline[linestyle=none](Q)(K)\aput[3pt]{0}{$4-x$}
		\ncline{A}{H}\bput[3pt]{0}{$2$}
		\ncline{B}{K}\bput[3pt]{0}{$1$}
		
		\multips(-1,0)(1,0){9}{
		\psline[linestyle=dotted,dotsep=1pt]{-}(0,0)(0,6)}
		\multips(-1,0)(0,1){7}{
		\psline[linestyle=dotted,dotsep=1pt]{-}(0,0)(8,0)}
		
		\psset{linecolor=blue,linewidth=1.3\pslinewidth}
		\pstLineAB{A}{P}
		\pstLineAB{Q}{P}
		\pstLineAB{Q}{B}
		
		\rput{0}(3,-1){Figura 5}

		\endpspicture
\end{center}


Possiamo ora minimizzare questa funzione ricorrendo al calcolo differenziale. Dallo studio del segno della derivata prima, cioè della disequazione (interessante di per sè)
$$\frac{x}{\sqrt{4+x^{2}}}>\frac{4-x}{\sqrt{1+(4-x)^{2}}},$$
segue che la funzione ha un minimo assoluto per $x=8/3$, e che tale minimo vale 6, come mostrato dal grafico qualitativo riportato qui sotto.
\begin{center}
	\psset{unit=2.5ex}
	\pspicture(-2,-1.5)(7,10)
	
	\psplot[linewidth=1.5\pslinewidth]{-2}{6}{x 2 exp 4 add sqrt 1 add 4 x sub 2 exp 1 add sqrt add}
	
	\psaxes[Dy=3,Dx=2,linewidth=0.7\pslinewidth]{->}(0,0)(-2,0)(7,10)
	\psline[linecolor=blue,linestyle=dashed]{-}(!8 3 div 0)(!8 3 div 6)(0,6)
	
	\endpspicture
\end{center}

La lunghezza minima è pertanto 600 metri.

Teoricamente l'approccio computazionale si può anche estendere al punto (b).  Come prima si osserva che ci si può limitare a considerare i percorsi costituiti da 3 tratti rettilinei e 2 ponti, uno sull'affluente ed uno sul fiume.  Tali percorsi sono univocamente determinati da 2 parametri $x$ e $y$: possiamo ad esempio prendere $x$ definita come nel caso precedente, e come $y$ la distanza in orizzontale tra $C$ e l'estremo piú vicino a $C$ del ponte sul fiume. In funzione di questi parametri si possono quindi calcolare le lunghezze dei 3 tratti del percorso: ci si riduce in questo modo a minimizzare la funzione di 2 variabili
$$\sqrt{4+x^{2}}+\sqrt{y^{2}+(10-y)^{2}}+\sqrt{y^{2}+(9-x-y)^{2}}.$$

Una possibile strategia è ora quella di pensare $y$ come un parametro fisso e minimizzare in $x$ la somma del primo e del terzo addendo con la stessa tecnica del passo precedente.  In questo modo si ottiene una funzione della sola $y$ da minimizzare a sua volta.  È chiaro che gli ingenti calcoli rendono questo approccio decisamente
svantaggioso in questo caso.}

\renewcommand{\Osservazioni}{
}

\renewcommand{\Commenti}{
}